% !PRINTED 645
\weberpage{645}

Ist aber $m$ eine Potenz von 2, so ist die Congruenz
\[
p^h \equiv -1 \ (\mathrm{mod}\ m)
\]
nur dann möglich, wenn $p \equiv -1\ (\mathrm{mod}\ m)$ ist. Denn jede gerade Potenz einer ungeraden Zahl ist $\equiv + 1\ (\mathrm{mod}\ 8)$ und kann also nicht $\equiv -1\ (\mathrm{mod}\ m)$ sein. Ist aber $h$ ungerade, so ist
\[
\frac{p^h + 1}{p + 1} = p^{h-1} - p^{h-2} + \cdots + 1
\]
selbst ungerade, und folglich ist $p^h + 1$ durch keine höhere Potenz von 2 theilbar, als $p + 1$. Demnach können wir dem Satze 2. noch folgende nähere Bestimmung hinzufügen:

3. \gesperrt{Wenn $m$ ungerade ist, so gehören die Primzahlen $p_1$ zu einem geraden, die Primzahlen $p_2$ zu einem ungeraden Exponenten.}

\gesperrt{Ist $m$ gerade, so gehören die Primzahlen der Form $km - 1$ zur ersten und alle anderen ungeraden Primzahlen zur zweiten Art.}

\begin{center}
\S.\,178.\\[4pt]
\gesperrt{Die Einheiten des Körpers $H_m$.}
\end{center}

Die dem Körper $H_m$ angehörigen Einheiten sind von besonderer Wichtigkeit für die Anwendungen. Auf sie lassen sich auch die Einheiten des Körpers $\Omega_m$ zurückführen nach folgendem Satze:

1. \gesperrt{Jede Einheit des Körpers $\Omega_m$ ist das Product einer Einheitswurzel und einer reellen Einheit des Körpers $\Omega_m$.}

Bezeichnen wir irgend eine Einheit des Körpers $\Omega_m$ mit $\mathfrak{E}(r)$, so ist $\mathfrak{E}(r^{-1})$ der conjugirt imaginäre Werth, der gleichfalls eine Einheit darstellt, und der Quotient $\mathfrak{E}(r) : \mathfrak{E}(r^{-1})$ ist wieder eine Einheit, deren absoluter Werth
\weq{1}{\sqrt{\dfrac{\mathfrak{E}(r)}{\mathfrak{E}(r^{-1})}\ \dfrac{\mathfrak{E}(r^{-1})}{\mathfrak{E}(r)}}}
gleich 1 ist, und folglich ist dieser Quotient eine Einheitswurzel und mithin $= \pm r^h$, worin $h$ irgend ein ganzzahliger Exponent ist (§. 175, Satz 1., 2.). Wir haben also
\weq{2}{\mathfrak{E}(r) = \pm r^h \mathfrak{E}(r^{-1}).}
